\chapter{The Arora--Safra Route: Polynomials, Sum-Check and Composition}
\label{lec:as}

{\sffamily\large\color{softgray}How a global algebraic statement becomes a handful of local queries\par}
\vspace{0.5em}

\noindent
Lectures~\ref{lec:why} and~\ref{lec:games} took the PCP theorem as a black box and asked what it gives us.
This lecture opens the box---or rather, opens the first one, the proof of \citet{AroraSafra98}.  The second
proof, the one that reaches $O(1)$ queries \citep{ALMSS98}, is Lecture~4; the third, combinatorial one is
Lecture~6.

Be warned about what Arora and Safra actually proved.  Their theorem is \emph{not} $\NP=\PCP[O(\log n),O(1)]$.  It is
\[
  \NP=\PCP[\log n,\log n]\qquad\text{and}\qquad
  \NP=\PCP\bigl[\log n,\ \log^{1/2+\eps}n\bigr]\ \text{ for every fixed }\eps>0,
\]
a verifier that reads a \emph{sublogarithmic} number of proof bits.  The step from $\log^{1/2+\eps}n$ to $O(1)$ was taken
soon afterwards by \citet{ALMSS98} (both papers appeared at FOCS 1992), using the same composition lemma and a sharper low-degree test (Section~\ref{sec:as-limits}).
We therefore state precisely the theorem proved here, and we resist the common shorthand ``Arora--Safra proved the PCP theorem''.

What makes the proof worth a full lecture is that it is built from \emph{modules} with clean interfaces, and we can
state and prove most of them.  The cast:
\begin{itemize}
  \item a \emph{code}: low-degree polynomials, which turn ``$f$ is almost right everywhere'' into ``$f$ is exactly right at a random point'';
  \item an \emph{arithmetization}: satisfiability of a $3$-CNF becomes ``a certain polynomial sums to $0$ over a cube'';
  \item the \emph{sum-check}: a table that lets a verifier check such a sum by reading $O(m)$ entries;
  \item the \emph{low-degree test}: a table that lets a verifier check that a function is close to a polynomial by reading $O(m)$ entries;
  \item \emph{composition}: the lemma that lets a verifier with many queries to long symbols be rewritten as one that reads
        fewer, shorter symbols.
\end{itemize}
Each module has an input, an output and a parameter change.  The point of the lecture is to track them.

\begin{goalbox}{After this lecture you should be able to\ldots}
\begin{enumerate}[leftmargin=*,label=(\arabic*)]
  \item state exactly what \citet{AroraSafra98} proved and which constants separate it from the PCP theorem (Section~\ref{sec:as-land});
  \item prove the Schwartz--Zippel lemma, and use it to decode a function that is close to a low-degree polynomial
        (Section~\ref{sec:poly});
  \item arithmetize a $3$-CNF as a vanishing condition on a cube, and compute the single nonzero point of $g_A$ on a
        small example (Section~\ref{sec:arith});
  \item prove the soundness $1-dl/q$ of the sum-check by induction and compute the cheating probability of a given table
        (Section~\ref{sec:sumcheck});
  \item prove that a function all of whose axis-parallel lines are low-degree is a polynomial, and explain why the real test
        needs a \emph{robust} version of this fact (Section~\ref{sec:ldt});
  \item prove the composition lemma, and carry out the parameter chase that ends at $\log^{1/2+\eps}n$ queried bits
        (Sections~\ref{sec:compose}--\ref{sec:chase}).
\end{enumerate}
\end{goalbox}


%% ===========================================================================
\section{What was proved, and when}
\label{sec:as-land}
%% ===========================================================================

Recall the notation of Definition~\ref{def:pcp}: $\PCP[r,q]$ is the class of languages with a non-adaptive verifier using
$O(r)$ random bits and $O(q)$ queries, perfect completeness and soundness~$\tfrac12$.  Trivially
$\PCP[r,q]\subseteq\mathsf{NTIME}\bigl(2^{O(r)}q+\poly(n)\bigr)$, so $\PCP[O(\log n),q]\subseteq\NP$ for every $q=\poly(n)$ (we used this in
Section~\ref{sec:pcp}).  The whole question is how \emph{small} $q$ can be.

\begin{table}[t]
\centering\small
\renewcommand{\arraystretch}{1.3}
\begin{tabular}{@{}>{\raggedright\arraybackslash}p{3.5cm}
                   >{\raggedright\arraybackslash}p{4.4cm}
                   >{\raggedright\arraybackslash}p{4.4cm}@{}}
\toprule
Result & Randomness and queries & Consequence for \sat{Max-Clique}\\
\midrule
$\NP=\PCP[0,\poly(n)]$ (definition) & none, all of the witness & none\\
\citet{BFLS91} & $\NP\subseteq\PCP[\polylog n,\polylog n]$ (implicit) & ---\\
\citet{FGLSS96} & $\PCP[\log n\log\log n,$\newline $\log n\log\log n]$ & hard within any constant unless $\NP\subseteq\mathsf{DTIME}(n^{O(\log\log n)})$\\
\citet{AroraSafra98} & $\PCP[\log n,\log n]$ and\newline $\PCP[\log n,\log^{1/2+\eps}n]$ & within $2^{c\log^{1/2-\eps}N}$ is $\NP$-hard\\
\citet{ALMSS98} & $\PCP[\log n,1]$ & within $N^{\eps_1}$ is $\NP$-hard\\
\bottomrule
\end{tabular}
\caption{The race for the number of queries (all verifiers non-adaptive; $N$ is the number of vertices of the clique instance).
The row of Arora and Safra is the subject of this lecture.  Each row implies the ones above it, since $\PCP[r,q]\subseteq\PCP[r',q']$ whenever $r\le r'$ and $q\le q'$.}
\label{tab:race}
\end{table}

\noindent
Two further statements from \citet{AroraSafra98} will be used.

\begin{theorem}[Arora--Safra]\label{thm:as}
For every fixed $\eps>0$, $\NP=\PCP\bigl[\log n,\ \log^{1/2+\eps}n\bigr]$.
\end{theorem}

\begin{corollary}[Arora--Safra]\label{cor:as-clique}
For all constants $c,\eps>0$, if \sat{Max-Clique} on $N$-vertex graphs can be approximated within $2^{c\log^{1/2-\eps}N}$
in polynomial time, then $\Ptime=\NP$.
\end{corollary}

\noindent
Corollary~\ref{cor:as-clique} is the \emph{reason} the paper was written: it removes the ``almost'' from the title of \citet{FGLSS96}, who had
shown that approximating Clique within any constant factor implies $\NP\subseteq\mathsf{DTIME}(n^{O(\log\log n)})$.  We derive it in Section~\ref{sec:clique}.

\begin{pitfall}{three different statements share the name ``PCP theorem''}
(1) $\NP=\PCP[\log n,\log n]$ and $\NP=\PCP[\log n,\log^{1/2+\eps}n]$: Arora--Safra, FOCS 1992.
(2) $\NP=\PCP[\log n,1]$: Arora--Lund--Motwani--Sudan--Szegedy, FOCS 1992.  This is what Lecture~\ref{lec:why} called the PCP theorem
(Theorem~\ref{thm:pcp1}); the \emph{approximation} consequences for \sat{Max-3SAT} need~(2), not~(1): through the reduction of Proposition~\ref{prop:pcp-equiv} the constant $\eps_0$ is $2^{-O(q)}$, so with
$\log^{1/2+\eps}n$ queried bits one only gets a gap that shrinks with~$n$, not a constant (Arora and Safra, Section~6).
(3) The combinatorial proof of the same statement as~(2) by \citet{Dinur07}.
Arora and Safra's own Section~6 is explicit about this division of labour.
\end{pitfall}


%% ===========================================================================
\section{Polynomials as codes}
\label{sec:poly}
%% ===========================================================================

Fix a finite field $\F_q$.  For integers $h,m\ge1$ let
\[
  \mathcal P^m_h\;=\;\{\,p\in\F_q[x_1,\dots,x_m]\ :\ \deg_{x_i}p\le h\text{ for every }i\,\}
\]
be the polynomials of degree at most $h$ \emph{in each variable}, viewed as functions $\F_q^m\to\F_q$.  (Each $p\in\mathcal P^m_h$ has total degree at most $hm$.)
For $f,g\colon\F_q^m\to\F_q$ write $\Delta(f,g)=\Pr_{x\in\F_q^m}[f(x)\ne g(x)]$ for their (normalized) distance, and say that $f$ is
\defn{$\delta$-close} to $\mathcal P^m_h$ if $\Delta(f,p)\le\delta$ for some $p\in\mathcal P^m_h$.

\begin{lemma}[Schwartz--Zippel]\label{lem:sz}
A nonzero $p\in\mathcal P^m_d$ vanishes on at most a $dm/q$ fraction of the points of $\F_q^m$.
\end{lemma}
\begin{proof}
Induction on $m$.  For $m=1$ a nonzero univariate polynomial of degree $\le d$ has at most $d$ roots.
For $m>1$ write $p=\sum_{i=0}^{k}x_1^i\,p_i(x_2,\dots,x_m)$ with $k\le d$ and $p_k\ne0$.  By induction $p_k$ vanishes on at most a
$d(m-1)/q$ fraction of $\F_q^{m-1}$.  For every point $(x_2,\dots,x_m)$ where $p_k\ne0$, the polynomial $p(\cdot,x_2,\dots,x_m)$ is a
nonzero univariate polynomial of degree $k\le d$, hence vanishes for at most $k$ values of $x_1$.  So the fraction of zeros is at most
$\frac{d(m-1)}q+\frac kq\le\frac{dm}q$.
\end{proof}

\begin{corollary}[Codes of large distance]\label{cor:dist}
Distinct $p,p'\in\mathcal P^m_h$ satisfy $\Delta(p,p')\ge1-hm/q$.  If $hm<q/2$ and $f$ is $\delta$-close to $\mathcal P^m_h$ for some $\delta<\tfrac14$,
then $f$ is $\delta$-close to a \emph{unique} $\tilde f\in\mathcal P^m_h$, and no other polynomial of $\mathcal P_h^m$ is within $\tfrac12-\delta$ of~$f$.
\end{corollary}
\begin{proof}
The first claim is Lemma~\ref{lem:sz} applied to $p-p'$.  If $p,p'$ were both within $\delta$ of $f$, the triangle inequality would give
$\Delta(p,p')\le2\delta<\tfrac12<1-hm/q$, a contradiction.  The last claim is the same computation with $\tfrac12-\delta$.
\end{proof}

\noindent
Corollary~\ref{cor:dist} is the reason polynomials are used: a function that is $1\%$ wrong still \emph{determines} the polynomial it is trying to be.
Next, how to turn a string into a polynomial.

\begin{lemma}[Polynomial extension]\label{lem:ext}
Let $h<q$, identify $\{0,1,\dots,h\}$ with the first $h+1$ elements of $\F_q$, and let $s\colon\{0,\dots,h\}^m\to\F_q$ be any function.  Then there is a \emph{unique}
$\hat s\in\mathcal P^m_h$ with $\hat s=s$ on $\{0,\dots,h\}^m$, namely
\[
  \hat s(x)\;=\;\sum_{u\in\{0,\dots,h\}^m}s(u)\,\Lambda_u(x),\qquad
  \Lambda_u(x)=\prod_{i=1}^m\lambda_{u_i}(x_i),\qquad
  \lambda_a(x)=\prod_{b\in\{0..h\},\,b\ne a}\frac{x-b}{a-b}.
\]
\end{lemma}
\begin{proof}
Each $\lambda_a$ has degree $h$, equals $1$ at $a$ and $0$ at the other points of $\{0..h\}$; hence $\Lambda_u$ is $1$ at $u$ and $0$
elsewhere on the cube, and $\hat s\in\mathcal P^m_h$ agrees with $s$ on the cube.  For uniqueness it suffices to show that $p\in\mathcal P^m_h$ vanishing on the cube is
zero.  Write $p=\sum_{j\le h}x_m^jp_j(x_1,\dots,x_{m-1})$.  For every fixed $x'\in\{0..h\}^{m-1}$ the univariate polynomial $p(x',\cdot)$ has degree $\le h$ and $h+1$ roots,
so it is zero, \ie $p_j(x')=0$ for all $j$.  Hence each $p_j$ vanishes on $\{0..h\}^{m-1}$ and by induction $p_j=0$.
\end{proof}

\begin{example}[The multilinear extension of XOR]\label{ex:xor}
For $m=2$, $h=1$ and $s(0,0)=s(1,1)=0$, $s(1,0)=s(0,1)=1$, Lemma~\ref{lem:ext} gives
\[
  \hat s(x_1,x_2)=x_1(1-x_2)+(1-x_1)x_2=x_1+x_2-2x_1x_2 .
\]
Over $\F_7$ it takes the value $3+5-30\equiv6$ at $(3,5)$.  The function $s$ is determined by $4$ values; its extension is a table with $q^2=49$ entries, and two distinct such extensions differ on at least $1-2/7=5/7$ of the $49$ points
(Corollary~\ref{cor:dist}), so even a corrupted table pins down $s$.
\end{example}

\noindent
\textbf{Encoding strings.}  Let $n=(h+1)^m$.  Identify $\{1,\dots,n\}$ with the cube $\{0,\dots,h\}^m$ (base-$(h+1)$ digits) and encode $a\in\{0,1\}^n\subset\F_q^n$ by
$\hat a\in\mathcal P^m_h$.  The encoding is injective, and the codeword is a table of $q^m$ field elements; if $q=\poly(h)$ then
$q^m=h^{O(m)}=\poly(n)$ since $(h+1)^m=n$.  This is the sense in which the encoding is ``efficient''.

\begin{pitfall}{a table that is close to a polynomial is not a polynomial}
The verifier will not hold a polynomial; it holds a \emph{table} $f\colon\F_q^m\to\F_q$ supplied by an adversary.  The low-degree test (Section~\ref{sec:ldt}) can only certify that
$f$ is $\delta$-close to some $\tilde f\in\mathcal P^m_h$ for a small constant~$\delta$.  Any later step that needs values of $\tilde f$ at a point $x$ reads $f(x)$ and is correct unless $x$
lies in the $\delta$-fraction where $f\ne\tilde f$.  The design principle of the whole verifier is therefore: \emph{never evaluate $\tilde f$ at a point the adversary can predict.}
Every point at which $f$ is read must be uniformly distributed, so that the error probability is $\le\delta$ per read.
\end{pitfall}


%% ===========================================================================
\section{Arithmetization: satisfiability as a vanishing condition}
\label{sec:arith}
%% ===========================================================================

Let $\varphi$ be a $3$-CNF with $n=(h+1)^m$ variables and $n$ clauses (pad with dummy variables and clauses).  Let $\mathbb K=\{0,1,\dots,h\}\subset\F_q$.  Identify variables and clauses with
points of the cube $\mathbb K^m$.  Every clause $c\in\mathbb K^m$ has an ordered triple of variables $v_1(c),v_2(c),v_3(c)\in\mathbb K^m$ and signs
$\varsigma_j(c)\in\{0,1\}$ ($\varsigma_j(c)=1$ if the $j$th literal is unnegated).  Define
\[
  \chi_j(c,v)=\begin{cases}1&\text{if }v=v_j(c)\\0&\text{otherwise}\end{cases}\quad\bigl(\mathbb K^m\times\mathbb K^m\to\{0,1\}\bigr),\qquad
  \varsigma_j\colon\mathbb K^m\to\{0,1\},
\]
and let $\hat\chi_j\in\mathcal P^{2m}_h$, $\hat\varsigma_j\in\mathcal P^m_h$ be their polynomial extensions (Lemma~\ref{lem:ext}); they depend only on
$\varphi$ and are computable in time $\poly(n,q^m)$.

For $A\in\mathcal P^m_h$, which we think of as the extension of a Boolean assignment, define $g_A\colon\F_q^{4m}\to\F_q$ by
\[
  g_A(z,w_1,w_2,w_3)\;=\;\prod_{j=1}^3\hat\chi_j(z,w_j)\cdot\bigl(\hat\varsigma_j(z)-A(w_j)\bigr),\qquad z,w_1,w_2,w_3\in\F_q^m .
\]

\begin{lemma}[Algebraic view of $3$SAT]\label{lem:arith}
Let $A\in\mathcal P^m_h$.
\begin{enumerate}[label=(\alph*)]
  \item If $a\in\{0,1\}^n$ satisfies $\varphi$ and $A=\hat a$, then $g_A(y)=0$ for every $y\in\mathbb K^{4m}=\{0,\dots,h\}^{4m}$.  If $a\in\{0,1\}^n$ does not satisfy $\varphi$, then $g_{\hat a}$
        is nonzero at some point of $\mathbb K^{4m}$.
  \item Conversely, for \emph{every} $A$: if $g_A$ vanishes on $\mathbb K^{4m}$ then $\varphi$ is satisfiable.  (The assignment $a'_v=\mathbf 1[A(v)=1]$ satisfies it.)
  \item $g_A\in\mathcal P^{4m}_{6h}$, and evaluating $g_A$ at a point $(z,w_1,w_2,w_3)$ requires only the values of $A$ at the three points
        $w_1,w_2,w_3\in\F_q^m$ (plus $\poly(n)$ work independent of $A$).  If the point is uniform in $\F_q^{4m}$ then each $w_j$ is uniform in $\F_q^m$.
\end{enumerate}
\end{lemma}
\begin{proof}
Fix $(c,v_1,v_2,v_3)\in\mathbb K^{4m}$.  On the cube the extensions $\hat\chi_j,\hat\varsigma_j$ equal $\chi_j,\varsigma_j$.  If $(v_1,v_2,v_3)\ne(v_1(c),v_2(c),v_3(c))$ then some $\chi_j(c,v_j)=0$ and $g_A=0$
at this point.  Otherwise
$g_A(c,v_1,v_2,v_3)=\prod_{j=1}^3\bigl(\varsigma_j(c)-A(v_j)\bigr)$, which vanishes iff $A(v_j)=\varsigma_j(c)$ for some $j$.

(a) For Boolean $a$ and $A=\hat a$ we have $A(v_j)=a_{v_j}$, and $a_{v_j(c)}=\varsigma_j(c)$ says exactly that the $j$th literal of $c$ is true.  So $g_A$ vanishes at $(c,v_1(c),v_2(c),v_3(c))$ iff $a$ satisfies $c$.

(b) If $A(v_j(c))=\varsigma_j(c)\in\{0,1\}$ then $a'_{v_j(c)}=\varsigma_j(c)$ (for $\varsigma_j=1$ because $A=1$; for $\varsigma_j=0$ because $A=0\ne1$).  So every clause has a literal made true by $a'$.

(c) Each factor is a product of two polynomials of degree $\le h$ in every variable; the three factors share~$z$, so the degree is $\le6h$ in each $z_i$, and $\le2h$ in each $w_{j,i}$.  The rest is read off
the formula; the $w_j$ are disjoint blocks of coordinates of a uniform point, hence uniform.
\end{proof}

\begin{pitfall}{the arithmetization certifies the \emph{rounded} assignment}
Part~(b) says that a table $A$ with $g_A\equiv0$ on the cube is satisfiable-by-rounding, not that its values on the cube are bits.  Indeed take any satisfying $a$ and change $a_v$ to $5$ for a variable $v$ that
occurs in clauses whose other literals are true: $g_A$ still vanishes on the cube.  This does no harm for the \emph{soundness} of a verifier for $\sat{3SAT}$ (an unsatisfiable $\varphi$ is rejected by every $A$, by (b)),
but it matters in Section~\ref{sec:normal}, where the verifier will be asked about the \emph{decoded string} of~$A$.  There one needs the further vanishing condition $A(A-1)\equiv0$ on the cube, and one more sum-check.
\citet{AroraSafra98} leave this implicit; the cost is a constant factor.
\end{pitfall}

\begin{example}[A formula, an assignment, one nonzero point]\label{ex:arith}
Take $h=1$, $m=2$ (so $n=4$ variables $x_0,\dots,x_3$, named by their binary index, and $4$ clauses) and
\[
  C_0=x_3\vee\neg x_2\vee\neg x_1,\quad C_1=\neg x_3\vee x_2\vee x_0,\quad C_2=x_1\vee x_0\vee\neg x_2,\quad C_3=\neg x_0\vee\neg x_2\vee\neg x_3 .
\]
The assignment $a=(x_0,x_1,x_2,x_3)=(0,0,0,1)$ falsifies exactly $C_1$ (here $x_3=1$, $x_2=0$, $x_0=0$).  Computing $g_{\hat a}$ at all $2^8=256$ points of $\{0,1\}^{8}$
over $\F_{31}$ gives a single nonzero value:
\[
  g_{\hat a}\bigl(\underbrace{(0,1)}_{C_1},\underbrace{(1,1)}_{x_3},\underbrace{(1,0)}_{x_2},\underbrace{(0,0)}_{x_0}\bigr)
  \;=\;(0-1)(1-0)(1-0)\;=\;-1,
\]
the product of $\varsigma_j-a_{v_j}$ for the three false literals.  So a single violated clause becomes a single nonzero point of a $256$-point cube:
a needle in a haystack.  We cannot hope to find it by sampling, and this is exactly why we multiply by the random \emph{zero-testers} of the next lemma, which turn
``somewhere nonzero'' into ``the weighted sum is nonzero''.
\end{example}

\begin{lemma}[Zero-testers]\label{lem:zt}
For $t=(t_1,\dots,t_{k})\in\F_q^{k}$ let $R_t\in\mathcal P^{k}_h$ be the polynomial $R_t(y)=\prod_{i=1}^{k}I_{t_i}(y_i)$,
where $I_s\in\mathcal P^1_h$ is the polynomial with $I_s(a)=s^a$ for $a=0,\dots,h$ (Lemma~\ref{lem:ext}).  Then for every $f\colon\{0..h\}^{k}\to\F_q$ that is not identically zero,
\[
  \Pr_{t\in\F_q^{k}}\Bigl[\sum_{y\in\{0..h\}^{k}}R_t(y)f(y)=0\Bigr]\ \le\ \frac{hk}{q}.
\]
In particular this is less than $\tfrac18$ when $q>8hk$ ($q>32hm$ for $k=4m$).
\end{lemma}
\begin{proof}
For $y\in\{0..h\}^{k}$ we have $R_t(y)=\prod_it_i^{y_i}$.  Hence $F(t):=\sum_yR_t(y)f(y)=\sum_yf(y)\prod_it_i^{y_i}$ is a polynomial in $t$ of degree $\le h$ in every variable, and it is
nonzero because distinct $y$ give distinct monomials and some $f(y)\ne0$.  By Lemma~\ref{lem:sz}, $\Pr_t[F(t)=0]\le hk/q$.
\end{proof}

\begin{example}[Zero-testers on Example~\ref{ex:arith}]\label{ex:zt}
Here $k=4m=8$ and $g_{\hat a}|_{\mathbb K^{4m}}$ has a single nonzero point, $y=(0,1\,|\,1,1\,|\,1,0\,|\,0,0)$ with value $-1$.  So $F(t)=-t_2t_3t_4t_5$ (the coordinates where $y$ has a~$1$), which vanishes iff one of these four
coordinates of $t$ is $0$:
\[
  \Pr_t[F(t)=0]\;=\;1-\bigl(\tfrac{30}{31}\bigr)^4\;\approx\;0.123\ \ \le\ \frac{hk}q=\frac8{31}\approx0.258 .
\]
(A simulation with $2\cdot10^5$ samples gives $0.123$.)  The field is too small for the hypothesis $q>8hk=64$ of the lemma, but the bound $hk/q$ itself holds.  Note how sparse $F$ is: a single violated clause
is a one-monomial polynomial, and the zero-tester detects it unless $t$ is unlucky in one of four coordinates.
\end{example}

\begin{corollary}[The summed polynomials $P_t$]\label{cor:P}
For $t\in\F_q^{4m}$ and $A\in\mathcal P^m_h$ put $P^A_t=R_t\cdot g_A\in\mathcal P^{4m}_{7h}$ ($R_t$ the zero-testers with $k=4m$, $q>32hm$).  Then
\begin{itemize}
  \item if $A=\hat a$ for a satisfying assignment $a$, then $\sum_{y\in\mathbb K^{4m}}P^A_t(y)=0$ for \emph{every} $t$;
  \item if $\varphi$ is \emph{unsatisfiable}, then for \emph{every} $A\in\mathcal P^m_h$ we have $\sum_{y\in\mathbb K^{4m}}P^A_t(y)\ne0$ for all but at most an $\tfrac18$ fraction of the $t$.
\end{itemize}
Moreover $P^A_t(y)$ at a uniformly random $y\in\F_q^{4m}$ is computed from three uniformly distributed values of $A$.
\end{corollary}
\begin{proof}
First bullet: Lemma~\ref{lem:arith}(a).  Second bullet: by Lemma~\ref{lem:arith}(b), $g_A\not\equiv0$ on $\mathbb K^{4m}$ when $\varphi$ is unsatisfiable; apply Lemma~\ref{lem:zt} to $f=g_A|_{\mathbb K^{4m}}$.
\end{proof}

\begin{bigpicture}{What arithmetization buys}
A $3$-CNF is a list of $n$ constraints, and a single violated clause is a needle in a haystack that a verifier reading $O(1)$ clauses cannot find.  After arithmetization the same fact is one equation,
``$\sum_{y\in\mathbb K^{4m}}P^A_t(y)=0$'', about a single low-degree polynomial for a random~$t$.  The sum has $(h+1)^{4m}=n^4$ terms, but the polynomial has structure---low degree in every variable---and
the \emph{sum-check} of the next section exploits that structure to verify the sum by evaluating the polynomial at \emph{one} random point.
\end{bigpicture}


%% ===========================================================================
\section{The sum-check protocol, as a table}
\label{sec:sumcheck}
%% ===========================================================================

The sum-check is due to \citet{LFKN92}; we give the non-adaptive, table-based form of \citet{AroraSafra98}.

Let $B\in\mathcal P^{\ell}_d$ and $H\subseteq\F_q$.  For $1\le i\le\ell$ and $a_1,\dots,a_{i-1}\in\F_q$ put
\[
  B_{a_1,\dots,a_{i-1}}(x_i)\;=\;\sum_{x_{i+1},\dots,x_\ell\in H}B(a_1,\dots,a_{i-1},x_i,x_{i+1},\dots,x_\ell),
\]
a univariate polynomial of degree $\le d$ in $x_i$ (for $i=1$ we write $B_\varnothing(x_1)$).  A \defn{table of partial sums} is any assignment $g$ of a univariate polynomial
$g_{a_1,\dots,a_{i-1}}$ of degree $\le d$ to every prefix $(a_1,\dots,a_{i-1})\in\F_q^{i-1}$, $1\le i\le\ell$; it has $\sum_i q^{i-1}\le q^\ell$ entries, each of $O(d\log q)$ bits.
The \emph{honest} table is $g_{a_1..a_{i-1}}=B_{a_1..a_{i-1}}$.

\begin{definition}[Sum-check test]\label{def:sumcheck}
Input: $B$ (available through oracle access), $H$, a claimed value $c\in\F_q$ and a table $g$.
Choose $a_1,\dots,a_\ell\in\F_q$ independently and uniformly; then accept iff
\begin{align*}
  &\sum_{x\in H}g_\varnothing(x)=c,\\
  &\sum_{x\in H}g_{a_1,\dots,a_{i-1}}(x)=g_{a_1,\dots,a_{i-2}}(a_{i-1})\quad(2\le i\le\ell),\\
  &g_{a_1,\dots,a_{\ell-1}}(a_\ell)=B(a_1,\dots,a_\ell).
\end{align*}
\end{definition}

\noindent
The test reads $\ell$ table entries and one value of $B$ at the uniform point $a=(a_1,\dots,a_\ell)$.  The entries to be read depend only on $a$, not on
earlier answers: the test is non-adaptive.  It uses $\ell\log q=\log|\F_q^\ell|$ random bits.

\begin{lemma}[Sum-check]\label{lem:sumcheck}
Let $B\in\mathcal P^\ell_d$, $c\in\F_q$.
\begin{enumerate}[label=(\alph*)]
  \item If $\sum_{x\in H^\ell}B(x)=c$, the honest table is accepted with probability~$1$.
  \item If $\sum_{x\in H^\ell}B(x)\ne c$, then \emph{every} table is accepted with probability at most $d\ell/q$.
\end{enumerate}
\end{lemma}
\begin{proof}
(a) Each check is an identity between partial sums.

(b) Induction on $\ell$.  For $\ell=1$ the test accepts iff $\sum_Hg_\varnothing=c$ and $g_\varnothing(a_1)=B(a_1)$.  Since $\sum_HB\ne c$ we have $g_\varnothing\ne B$, and two distinct polynomials of
degree $\le d$ agree in at most $d$ points, so the acceptance probability is at most $d/q$.

For $\ell>1$ let $B_\varnothing(x_1)=\sum_{x_2..x_\ell\in H}B(x_1,x_2,\dots,x_\ell)$.  Then $\sum_{x\in H}B_\varnothing(x)=\sum_{H^\ell}B\ne c$.
\emph{Case 1:} $g_\varnothing=B_\varnothing$.  Then $\sum_Hg_\varnothing\ne c$ and the first check fails with probability~$1$.
\emph{Case 2:} $g_\varnothing\ne B_\varnothing$.  They agree in at most $d$ points $a_1$.  For every other $a_1$, the claim ``$\sum_{x_2..x_\ell\in H}B(a_1,x_2,\dots,x_\ell)=g_\varnothing(a_1)$'' is
false.  Fixing $a_1$, the remaining checks form exactly the sum-check test for the polynomial $B(a_1,\cdot)\in\mathcal P^{\ell-1}_d$, the claim $c'=g_\varnothing(a_1)$ and the sub-table $g_{a_1,\ast}$;
by induction it accepts with probability $\le d(\ell-1)/q$.  Hence the overall acceptance probability is at most
\[
  \frac dq\cdot1+\Bigl(1-\frac dq\Bigr)\frac{d(\ell-1)}q\ \le\ \frac{d\ell}q.\qedhere
\]
\end{proof}

\noindent
(The recursion $p_\ell\le d/q+(1-d/q)p_{\ell-1}$, $p_1\le d/q$, solves to the slightly sharper $p_\ell\le1-(1-d/q)^\ell\le d\ell/q$.)

\begin{example}[A cheating table, exactly]\label{ex:sc}
Let $q=7$, $d=2$, $\ell=2$, $H=\{0,1\}$ and $B(x_1,x_2)=x_1x_2+x_2^2+1$.  Then
$B_\varnothing(x_1)=B(x_1,0)+B(x_1,1)=1+(x_1+2)=x_1+3$, and the true sum is $B_\varnothing(0)+B_\varnothing(1)=3+4=7\equiv0$.
Suppose the prover wants to convince us of the false claim $c=1$.  An exhaustive computation over all tables (every polynomial of degree $\le2$ at every node of the tree) gives
\[
  \max_g\Pr[\text{accept}]\;=\;\frac{24}{49}\;=\;1-\Bigl(\frac57\Bigr)^2\;\approx\;0.490\ \ <\ \frac{d\ell}q=\frac47\approx0.571 .
\]
The maximum is the same for every false claim $c\ne0$, and equals exactly the sharper bound $1-(1-d/q)^\ell$ of the remark after the proof: at the first level the cheating polynomial can have
$d=2$ points $a_1$ at which it agrees with $B_\varnothing$, and at each of the other points the second level leaves a further $d/q$.  The bound $d\ell/q$ of the lemma is its linearization.
\end{example}

\begin{pitfall}{a table of partial sums is a prover strategy written out}
In the interactive sum-check of \citet{LFKN92} the prover answers round by round, and its $i$th message may depend on the challenges $a_1,\dots,a_{i-1}$.  A deterministic strategy of such a prover is
\emph{exactly} a table of partial sums: the entry $g_{a_1,\dots,a_{i-1}}$ is its $i$th message.  So the lemma says the same thing in both models.  What the table form adds is that the verifier
may choose all of $a$ first and then read the $\ell$ entries on the path to $a$ in one batch: the test is non-adaptive.  This is needed for the composition lemma (Section~\ref{sec:compose}).
\end{pitfall}

\subsection{Putting arithmetization and sum-check together}

Combine Corollary~\ref{cor:P} and Lemma~\ref{lem:sumcheck} with $B=P^A_t$, $d=7h$, $\ell=4m$, $H=\{0..h\}$, $c=0$.  A verifier that holds $f\colon\F_q^m\to\F_q$ (in the role of $A$)
and a collection of sum-check tables $\{g^{(t)}\}_{t\in\F_q^{4m}}$ does the following:
\begin{enumerate}
  \item pick $t\in\F_q^{4m}$ uniformly ($4m\log q$ random bits) and compute from $\varphi$ and $t$ everything about $P^f_t$ except the values of $f$;
  \item run the sum-check on $P^f_t$ with claim $c=0$ and table $g^{(t)}$.
\end{enumerate}
If $f$ extends a satisfying assignment, there are tables making this accept with probability~$1$.  If $\varphi$ is unsatisfiable and $f$ is \emph{any} polynomial in $\mathcal P^m_h$, then by
Corollary~\ref{cor:P} and Lemma~\ref{lem:sumcheck} the test accepts with probability at most $\tfrac18+28hm/q$ (here $d\ell/q=7h\cdot4m/q$).  In the next section we remove the assumption that $f$ is a polynomial.


%% ===========================================================================
\section{The low-degree test}
\label{sec:ldt}
%% ===========================================================================

The sum-check of the previous section was proved for polynomials.  The verifier holds a table $f$.  We need a \emph{local} test that accepts every
$f\in\mathcal P^m_h$ and rejects every $f$ far from $\mathcal P^m_h$.

\subsection{The exact statement, and why it is not enough}

An \defn{axis-parallel line} of direction $i$ in $\F_q^m$ is $\ell=\{(a_1,\dots,a_{i-1},x,a_{i+1},\dots,a_m):x\in\F_q\}$.  There are $m\,q^{m-1}$ of them.
If $p\in\mathcal P^m_h$ then its restriction to every axis-parallel line is a univariate polynomial of degree $\le h$.  The converse is true, and its proof is a good warm-up.

\begin{lemma}[Axis-parallel lines suffice, exact version]\label{lem:axis}
Let $f\colon\F_q^m\to\F_q$ be such that the restriction of $f$ to every axis-parallel line is a polynomial of degree $\le h$ (with $h<q$).  Then $f\in\mathcal P^m_h$.
\end{lemma}
\begin{proof}
Induction on $m$; $m=1$ is the hypothesis.  Fix distinct $c_0,\dots,c_h\in\F_q$.  For each $j$ the function $f(\cdot,c_j)$ on $\F_q^{m-1}$ satisfies the hypothesis in $m-1$ variables, so it is a polynomial
$p_j\in\mathcal P^{m-1}_h$.  Let $\lambda_j$ be the univariate Lagrange polynomials of degree $h$ for the nodes $c_0,\dots,c_h$ and define
$G(x',x_m)=\sum_{j=0}^h\lambda_j(x_m)\,p_j(x')\in\mathcal P^m_h$.  For each fixed $x'$, the functions $f(x',\cdot)$ and $G(x',\cdot)$ are both univariate polynomials of degree $\le h$ and they agree at the
$h+1$ points $c_0,\dots,c_h$; hence they are equal.  So $f=G$ everywhere.
\end{proof}

\noindent
As a sanity check, take $q=5$, $h=1$, $m=2$.  A pruned exhaustive search through the $5^{25}$ functions $\F_5^2\to\F_5$ finds exactly $625=5^{(h+1)^2}$ with all rows and all columns of degree $\le1$:
precisely the $5^4$ polynomials of $\mathcal P^2_1$, as the lemma predicts.

\begin{bigpicture}{Exact versus robust}
Lemma~\ref{lem:axis} is a statement about \emph{perfect} functions.  A verifier can only look at a few lines, and can at best find that \emph{most} of the lines it looked at are low-degree and \emph{mostly} agree where they meet.
The theorem the test needs is the \emph{robust} form: if the restrictions to most axis-parallel lines are consistent with each other and with~$f$, then $f$ is \emph{close} to a polynomial.
All the difficulty of the Arora--Safra analysis sits here, and it is the reason the field size is cubic in the degree.
\end{bigpicture}

\subsection{The test of Arora and Safra}

A \defn{decomposition} is a table $g$ that assigns to every axis-parallel line $\ell$ a univariate polynomial $g[\ell]$ of degree $\le h$ (an entry of $O(h\log q)$ bits).  For $1\le i\le m$ and
$a\in\F_q^m$ let $g_i(a)=g[\ell_i(a)](a_i)$, where $\ell_i(a)$ is the line of direction $i$ through $a$.  A \defn{checkpoint} is a pair $(i,a)$ with $1\le i<m$; the decomposition is \defn{consistent} at it if
$g_i(a)=g_{i+1}(a)$.  Note that a genuine decomposition of a polynomial $p$ ($g[\ell]=p|_\ell$) is consistent everywhere and $g_i=p$ for all $i$.

\begin{definition}[Low-degree test]\label{def:ldt}
Given $f\colon\F_q^m\to\F_q$, the closeness parameter $\delta<0.1$ and a decomposition $g$:
\begin{enumerate}[label=(\roman*)]
  \item choose $\lceil4/\delta\rceil$ points $a\in\F_q^m$ uniformly and independently, and check $f(a)=g_1(a)$ at each;
  \item choose $\lceil8m/\delta\rceil$ checkpoints $(i,a)$ uniformly and independently, and check $g_i(a)=g_{i+1}(a)$ at each.
\end{enumerate}
Accept iff all checks pass.
\end{definition}

\noindent
The test reads $O(1/\delta)$ values of $f$ and $O(m/\delta)$ entries of $g$: $O(m)$ queries for constant $\delta$, non-adaptively.  (The randomness is $O(m\log q)$ bits per sample, which is reduced to $O(\log q^m)$ in total
by pairwise-independent sampling; we omit it.)  It accepts every $f\in\mathcal P^m_h$ together with its true decomposition.

\begin{lemma}[Robust consistency; Arora--Safra, Lemma~4.6.4]\label{lem:robust}
Let $m\ge2$ and $q\ge100\,h^3m^2$.  If $g_1$ is not $\eps$-close to $\mathcal P^m_h$, then the fraction of checkpoints at which $g$ is inconsistent is at least
$\gamma^{2(m-1)}\min\{0.1,\eps\}/(m-1)$, where $\gamma=1-\sqrt{h/q}$.  Since $\sqrt{h/q}\le1/(10m)$ we have $\gamma^{2(m-1)}\ge1-2(m-1)/(10m)\ge0.8$, so this fraction is at least $\min\{0.1,\eps\}/(2m)$.
\end{lemma}

\noindent
We do not prove Lemma~\ref{lem:robust}.  The proof is an induction on $m$.  Its base case $m=2$ is the \emph{bivariate} lemma that \citet{ALMSS98} quote as Theorem~69 (due to Arora--Safra): if the row polynomial and the column polynomial of a
function $F\colon\F_q^2\to\F_q$ each agree with $F$ on all but an $\eta$ fraction of points ($\eta\le\eta_0$) and $q\ge c\,h^3$, then $F$ is $4\eta$-close to a polynomial of degree $\le h$ in each variable.  This is the robust analogue
of the exact Lemma~\ref{lem:axis} with $m=2$, and it is where the \emph{cubic} dependence on~$h$ comes from.  \citet{PolishchukSpielman94} later gave a proof that needs only $q=\Omega(h)$.

\begin{corollary}[Soundness of the test]\label{cor:ldt}
Let $q\ge100\,h^3m^2$ and $m\ge2$.  If $f$ is not $\delta$-close to $\mathcal P^m_h$ ($\delta<0.1$), then for \emph{every} decomposition $g$ the test of Definition~\ref{def:ldt} rejects with probability at least $3/4$.
\end{corollary}
\begin{proof}
If $\Delta(f,g_1)\ge\delta/2$, part~(i) fails to reject with probability at most $(1-\delta/2)^{4/\delta}\le e^{-2}\approx0.135<\tfrac14$.
Otherwise $\Delta(f,g_1)<\delta/2$, so $g_1$ is not $\delta/2$-close to $\mathcal P^m_h$ (else $f$ would be $\delta$-close by the triangle inequality).  By Lemma~\ref{lem:robust} with $\eps=\delta/2\le0.05$ a fraction at least $\delta/(4m)$ of the
checkpoints is inconsistent, and part~(ii) misses all of them with probability at most $(1-\delta/4m)^{8m/\delta}\le e^{-2}<\tfrac14$.
\end{proof}

\begin{remark}
With $\delta=0.01$ the two failure probabilities are $\le0.1347$ and $\le0.1353$ respectively (computed from the formulas above).  The older analysis of this test (for $h=1$ in \citealp{FGLSS96}) required $O(mh)$ queries; Arora and Safra's algebraic argument gets $O(m)$, at the price of $q=\Omega(h^3m^2)$.  The number of queries to the low-degree table is exactly what is later
reduced to a constant by \citet{ALMSS98}; see Section~\ref{sec:as-limits}.
\end{remark}


%% ===========================================================================
\section{The basic verifier}
\label{sec:vpoly}
%% ===========================================================================

We now assemble the three modules.  Let $n=(h+1)^m$ with $h\ge\log n$ and $q=\Theta(h^3m^2)$ a prime power (note that $m<h$ forces $q=O(h^5)$).  The verifier $V_{\rm poly}$ for a $3$-CNF $\varphi$ with $n$ variables expects the proof to consist of

\begin{center}\small
\renewcommand{\arraystretch}{1.25}
\begin{tabular}{@{}>{\raggedright\arraybackslash}p{2.6cm}
                   >{\raggedright\arraybackslash}p{6.0cm}
                   >{\raggedright\arraybackslash}p{4.0cm}@{}}
\toprule
component & content & size\\
\midrule
$f$ & a function $\F_q^m\to\F_q$ (the claimed extension of a satisfying assignment) & $q^m=n^{O(1)}$ field elements\\
$g$ & a decomposition for the low-degree test of $f$ & $mq^{m-1}$ polynomials of degree $h$\\
$\{\tau^{(t)}\}_{t\in\F_q^{4m}}$ & a table of partial sums for the sum-check of $P_t^{f}$, one per $t$ & $\le2q^{8m}=n^{O(1)}$ polynomials of degree $7h$\\
\bottomrule
\end{tabular}
\end{center}

\noindent
and works as follows.

\begin{definition}[The verifier $V_{\rm poly}$]
Given the proof components above, $V_{\rm poly}$ does the following.
\begin{enumerate}
  \item Run the low-degree test (Definition~\ref{def:ldt}) on $f$ and $g$ with $\delta=0.01$; reject if it rejects.
  \item Choose $t\in\F_q^{4m}$ uniformly and run the sum-check (Definition~\ref{def:sumcheck}) for $B=P^{f}_t$ (that is, with the values of $f$ in place of $A$), $H=\{0..h\}$, $\ell=4m$, $d=7h$, $c=0$, and table $\tau^{(t)}$;
        reject if it rejects.  Whenever the sum-check needs $f$ at a point $w$, it reads $f(w)$.
  \item Otherwise accept.
\end{enumerate}
\end{definition}

\begin{theorem}[The basic verifier; Arora--Safra, Theorem~3.5(a)]\label{thm:vpoly}
$V_{\rm poly}$ is a non-adaptive verifier for $\sat{3SAT}$ with perfect completeness, soundness $<\tfrac14$, using $O(\log n)$ random bits and $O(m)$ queries,
each to an entry of $O(h\log h)$ bits, and with decision time $\poly(h)$.
\end{theorem}
\begin{proof}
\emph{Completeness.}  If $\varphi$ is satisfiable, let $f$ be the extension of a satisfying assignment, $g$ its true decomposition, and $\tau^{(t)}$ the honest table for $P_t^f$.  The low-degree test accepts by completeness, and
the sum-check accepts by Lemma~\ref{lem:sumcheck}(a), because $\sum_{y\in\mathbb K^{4m}}P^f_t(y)=0$ by Corollary~\ref{cor:P}.

\emph{Soundness.}  Let $\varphi$ be unsatisfiable and fix any proof.  If $f$ is not $0.01$-close to $\mathcal P^m_h$, step~1 rejects with probability $\ge3/4$ (Corollary~\ref{cor:ldt}).
So assume $f$ is $0.01$-close, and let $\tilde f\in\mathcal P^m_h$ be its unique nearest polynomial (Corollary~\ref{cor:dist}).  Step~2 accepts only if one of three events happens:
\begin{enumerate}[label=(\roman*)]
  \item $\sum_{\mathbb K^{4m}}P^{\tilde f}_t=0$: by Corollary~\ref{cor:P} this has probability $\le\tfrac18$ over $t$ (this is where unsatisfiability is used, for the arbitrary polynomial $\tilde f$);
  \item the sum is nonzero but the sum-check does not detect it: probability $\le d\ell/q=28hm/q$ by Lemma~\ref{lem:sumcheck}(b), which is applied to the genuine polynomial $B=P_t^{\tilde f}$;
  \item the sum-check reads $f(w_j)$ at one of three points $w_1,w_2,w_3$ with $f(w_j)\ne\tilde f(w_j)$: each $w_j$ is uniform (Lemma~\ref{lem:arith}(c)), so by the union bound this has probability $\le3\cdot0.01$.
\end{enumerate}
Altogether the acceptance probability is at most $\tfrac18+0.03+28hm/q<\tfrac14$ (for $q\ge h^3m^2$ and $h\ge\log n$ the last term is $o(1)$).

\emph{Parameters.}  The random bits are those of $t$, the point $a\in\F_q^{4m}$ of the sum-check, and the samples of the low-degree test: $O(m\log q)=O(m\log h)=O(\log n)$ in total, since $(h+1)^{m}=n$.
The low-degree test reads $O(m)$ entries and the sum-check $4m$ entries of $\tau^{(t)}$, plus three values of $f$.  Entries have $O(h\log q)=O(h\log h)$ bits.  Before reading the proof, the verifier computes from $\varphi$ and $(t,a)$ everything about
$P^f_t(a)$ that does not involve $f$: $\hat\chi_j,\hat\varsigma_j,R_t$ at the relevant points, which takes $\poly(n)$ time; after reading, it does $\poly(m,h,\log q)=\poly(h)$ field operations.  All positions queried are determined by the random string alone: the verifier draws all its randomness first, so it can read the entries of the
low-degree test and of the sum-check in one batch.  This is the point of using sum-check \emph{tables} rather than the interactive protocol.
\end{proof}

\begin{remark}
Note how the verifier avoids a circularity.  The sum-check needs $\tilde f$ at three points, but only $f$ is available.  The three points are \emph{uniform}, hence outside the $1\%$ of disagreement with probability $\ge0.97$.
This is the ``never evaluate at a predictable point'' principle of Section~\ref{sec:poly}, and it is the reason the arithmetization of Lemma~\ref{lem:arith}(c) was designed so that each $w_j$ is uniform.
\end{remark}

Theorem~\ref{thm:vpoly} is already a PCP: choosing $h+1=n^{1/m}$ for a constant $m$ it has $O(1)$ queries.  The catch is the \emph{size} of the queried symbols: $O(h\log h)=n^{\Omega(1/m)}$ bits each.
To get a verifier that reads \emph{bits} one must re-encode every queried symbol and check it without reading it.  That is the job of composition.


%% ===========================================================================
\section{Normal form verifiers and split-encoded assignments}
\label{sec:normal}
%% ===========================================================================

Suppose an \emph{outer} verifier $V_1$ has a decision circuit $C_w$ (on the symbols it queried under random string $w$), of size $T$ much smaller than the proof.
Instead of reading the queried symbols $P[i_1(w)],\dots,P[i_Q(w)]$ in full, we would like to check ``$C_w$ accepts them'' by running a second, \emph{inner} verifier $V_2$ on the formula of $C_w$.
Two requirements:
\begin{itemize}
  \item the inner verifier must work when the \emph{input assignment is spread over several separate pieces}, each of which is held in some encoded form in the proof (the $Q$ queried symbols and a scratch string);
  \item the inner verifier must read only a few \emph{bits} of each piece, since reading the pieces in full is what we want to avoid.
\end{itemize}
This is formalized in the following definition (a simplified form of Definition~3.2 of \citet{AroraSafra98}).

\begin{definition}[Normal form verifier]\label{def:normal}
An \defn{$(r,s,q,t)$-constrained normal form verifier} $V_2$ has an associated encoding scheme $\sigma$ (bit strings to codewords of minimum distance $\delta_{\min}$ over an alphabet of $2^{O(s)}$ symbols).
Given a $3$-CNF $\psi$ with $N$ variables and an integer $p$ dividing $N$, it is given access to a proof $z_1\sharp\cdots\sharp z_p\sharp\pi$, where each $z_i$ is a word over the alphabet and $\pi$ is auxiliary, and must satisfy:
\begin{itemize}
  \item \emph{(completeness)} if each $z_i=\sigma(a_i)$ for $a_i\in\{0,1\}^{N/p}$ and $a_1\circ\dots\circ a_p$ satisfies $\psi$, there is a $\pi$ with acceptance probability~$1$;
  \item \emph{(soundness)} if some $z_i$ is not $\delta_{\min}/3$-close to a codeword, or if each $z_i$ is close to a codeword $\sigma(a_i)$ but $a_1\circ\dots\circ a_p$ does not satisfy $\psi$, then every $\pi$ is accepted with probability $<\tfrac12$;
  \item \emph{(resources)} it uses $O(r(N))$ random bits and alphabet symbols of $O(s(N))$ bits, makes $O(p+q(N))$ non-adaptive queries, and has decision time $O(p\,t(N))$.
\end{itemize}
\end{definition}

\noindent
The unusual feature is the query count $O(p+q(N))$ rather than $O(p\cdot q(N))$: the inner verifier reads $O(1)$ symbols of each piece plus $q(N)$ symbols of the \emph{shared} auxiliary table.
Proposition~3.3 of \citet{AroraSafra98} records the trivial case $p=1$: a normal form verifier gives a verifier for $\sat{3SAT}$ with the same resources.

This notion is the ancestor of the \emph{PCP of proximity} \citep{BGHSV06} and of the \emph{assignment tester} \citep{DR04}, which isolate exactly the property used here---``a tester that looks at a few places of an encoded input
and a proof''---and which are what \citet{Dinur07} composes with in Lecture~6.

\subsection{How $V_{\rm poly}$ becomes a normal form verifier}

By Theorem~\ref{thm:vpoly}, $V_{\rm poly}$ is a verifier with parameters $(\log n,\ h\log h,\ m,\ \poly(h))$ for $r,s,q,t$.  Two additions turn it into a normal form verifier, and
both are instances of the exact-versus-robust pattern.

\paragraph{The encoding.}  Take $\sigma$ to be the polynomial extension: $\sigma\colon\{0,1\}^N\to\mathcal P^{m}_h$ with $N=(h+1)^m$, a code of minimum distance $\ge1-hm/q>\tfrac12$.

\paragraph{Split encodings.}  Suppose $N=(h+1)^m$ and $p=(h+1)^{m-\ell}$, so that each of the $p$ pieces has $N/p=(h+1)^\ell$ bits.  The proof contains $f\colon\F_q^m\to\F_q$ (the encoding of the whole assignment, for the sum-check as before)
and $f_1,\dots,f_p\colon\F_q^\ell\to\F_q$ (the encodings of the pieces), and must check the \defn{concatenation property}: for $i\in\mathbb K^{m-\ell}$ and $u\in\mathbb K^\ell$,
\begin{equation}\label{eq:concat}
  f(i,u)=f_i(u).
\end{equation}

\begin{lemma}[Concatenation]\label{lem:concat}
Let $f\in\mathcal P^m_h$ and $f_1,\dots,f_p\in\mathcal P^\ell_h$ (indexed by $i\in\mathbb K^{m-\ell}$).  Then~\eqref{eq:concat} holds if and only if, as polynomials,
\[
  f(x_1,\dots,x_m)\;=\;\sum_{i\in\mathbb K^{m-\ell}}\Lambda_i(x_1,\dots,x_{m-\ell})\,f_i(x_{m-\ell+1},\dots,x_m),
\]
where $\Lambda_i$ is the Lagrange basis polynomial of Lemma~\ref{lem:ext} for the cube $\mathbb K^{m-\ell}$.
\end{lemma}
\begin{proof}
Let $G$ denote the right-hand side; $G\in\mathcal P^m_h$.  On a cube point $(i',u)$ we have $\Lambda_i(i')=\mathbf 1[i=i']$, so $G(i',u)=f_{i'}(u)$.  If~\eqref{eq:concat} holds, $f$ and $G$ agree on $\mathbb K^m$ and are equal by the
uniqueness in Lemma~\ref{lem:ext}.  Conversely, if $f=G$ then $f(i',u)=G(i',u)=f_{i'}(u)$.
\end{proof}

\noindent
So the verifier tests the identity $f(x)=\sum_i\Lambda_i(x')f_i(x'')$ at $O(1)$ random points $x=(x',x'')\in\F_q^m$.  Each test reads $f(x)$ and one value $f_i(x'')$ for each of the $p$ pieces: $O(1)$ values per piece, $O(p)$ in all,
which is the source of the $O(p+q)$ in Definition~\ref{def:normal}.  If the functions are polynomials but the identity fails, then $f-G$ is a nonzero polynomial and the test catches it with probability $\ge1-hm/q$ per sample (Lemma~\ref{lem:sz}).
Arora and Safra prove the \emph{robust} version (their Theorem~4.4.2): if $f$ is $0.01$-close and the test accepts with probability $>\tfrac14$, then each $f_i$ is $0.03$-close and the concatenation property holds for the nearest polynomials.
The proof rests on a second robust lemma (their Lemma~5.1.1, proved with the same field-size hypothesis as Lemma~\ref{lem:robust}): a function $G$ that is $0.02$-close to a polynomial and all of whose
restrictions $x'\mapsto G(x',x'')$, for fixed $x''$, are polynomials of degree $\le h$, has the property that every restriction $x''\mapsto G(x',x'')$ is $0.03$-close to a polynomial---namely to the corresponding
restriction of the nearest polynomial $\tilde G$.  Applied to $G=\sum_i\Lambda_if_i$ it yields that each $f_i$ is $0.03$-close, with $\tilde f_i=\tilde f(i,\cdot)$.

\begin{pitfall}{decoding must produce bits}
In the definition of a normal form verifier the decoded pieces $a_i=\sigma^{-1}(\text{nearest codeword to }z_i)$ must be \emph{bit strings}, and $\psi$ is evaluated on their concatenation.
By the pitfall after Lemma~\ref{lem:arith}, the arithmetization alone only certifies satisfiability of the rounded assignment; to meet the definition the verifier adds a Booleanity check
($A(A-1)$ vanishing on the cube) to the sum-check, which changes only the constants.
\end{pitfall}


%% ===========================================================================
\section{The composition lemma}
\label{sec:compose}
%% ===========================================================================

Let $V_1$ be a verifier with random bits $R$, alphabet of $2^{S}$ symbols, $Q$ queries, decision time $T$, perfect completeness and soundness $\tfrac14$.  For a random string $w\in\{0,1\}^R$ it queries positions
$i_1(w),\dots,i_Q(w)$ of its proof $P$, and accepts iff a circuit $C_w$ of size $O(T^2)$ accepts $P[i_1(w)]\circ\dots\circ P[i_Q(w)]$.  By the Cook--Levin theorem for circuits this is equivalent to the satisfiability of a $3$-CNF
$c_w$ of size $\tau=O(T^2)$ over the bits of the $Q$ symbols and of a witness string $y_w$.  We append the strings $y_w$ ($w\in\{0,1\}^R$) to $P$ as extra symbols and let $V_1$ read $y_w$ as a $(Q{+}1)$st query.
After this harmless change, writing $i_{Q+1}(w)$ for the position of $y_w$,
\begin{equation}\label{eq:decide}
  V_1\text{ accepts }P\text{ on }w\iff P[i_1(w)]\circ\dots\circ P[i_{Q+1}(w)]\text{ satisfies }c_w .
\end{equation}

\begin{lemma}[Composition; Arora--Safra, Lemma~3.4]\label{lem:compose}
Let $V_1$ be as above, with parameters $(R,S,Q,T)$, and let $V_2$ be an $(r,s,q,t)$-constrained normal form verifier (Definition~\ref{def:normal}).
Then the same language has a verifier $V_{\rm new}$ with perfect completeness and soundness $<\tfrac12$, and with
\[
  \begin{aligned}
  \text{random bits } &R+O(r(\tau)),\qquad &\text{alphabet bits } &O(s(\tau)),\\
  \text{queries } &O\bigl(Q+q(\tau)\bigr),\qquad &\text{decision time } &O\bigl(Q\,t(\tau)\bigr),
  \end{aligned}\qquad\tau=O(T^2).
\]
\end{lemma}

\begin{figure}[t]
\centering
\begin{tikzpicture}[x=1cm,y=1cm,every node/.style={font=\small},
   wd/.style={draw,minimum width=0.62cm,minimum height=0.62cm,inner sep=0pt,fill=white},
   hl/.style={wd,fill=accent!30},
   hy/.style={wd,fill=warn!35},
   hp/.style={wd,fill=claim!22}]
  \foreach \i in {0,...,13} \node[wd] at (\i*0.66,0) {};
  \node[hl] at (1*0.66,0) {};  \node[hl] at (5*0.66,0) {};  \node[hl] at (9*0.66,0) {};
  \node[hy] at (12*0.66,0) {};
  \foreach \i in {15,...,19} \node[wd] at (\i*0.66,0) {};
  \node[hp] at (17*0.66,0) {};
  \node[anchor=east] at (-0.5,0) {$P_{\rm new}$};
  \draw[decorate,decoration={brace,amplitude=3pt,mirror}] (-0.28,-0.5) -- (13*0.66+0.28,-0.5)
        node[midway,below=4pt,font=\footnotesize] {$Y$ words: $\sigma(P[j])$ and $\sigma(y_w)$};
  \draw[decorate,decoration={brace,amplitude=3pt,mirror}] (15*0.66-0.28,-0.5) -- (19*0.66+0.28,-0.5)
        node[midway,below=4pt,font=\footnotesize] {$2^R$ words $\pi_w$};
  % verifiers
  \node[draw,rounded corners=3pt,fill=black!5,minimum width=3.2cm,minimum height=0.8cm] (V1) at (3.0,2.4) {$V_1$: picks $w$, forms $c_w$};
  \node[draw,rounded corners=3pt,fill=black!5,minimum width=4.2cm,minimum height=0.8cm] (V2) at (9.6,2.4) {$V_2$ on $c_w$ (normal form)};
  \draw[-{Stealth[length=2mm]}] (V1)--(V2) node[midway,above,font=\footnotesize] {$c_w$};
  \foreach \x in {1,5,9,12} \draw[accent,-{Stealth[length=1.6mm]}] (V2.south) to[out=-110,in=90] (\x*0.66,0.35);
  \draw[claim,-{Stealth[length=1.6mm]}] (V2.south) to[out=-70,in=90] (17*0.66,0.35);
\end{tikzpicture}
\caption{The composed verifier $V_{\rm new}$ (Lemma~\ref{lem:compose}).  The proof consists of encodings $\sigma(P[j])$ of the symbols of $V_1$'s proof (including the witness strings $y_w$) and one auxiliary word
$\pi_w$ per random string $w$ of $V_1$.  On random string $w$ the verifier $V_1$ would read the $Q$ blue words and the orange word in full; $V_{\rm new}$ instead hands them to $V_2$ as a split-encoded assignment to $c_w$
and reads only a few symbols of each.}
\label{fig:compose}
\end{figure}

\begin{proof}
Replacing $V_1$ and $V_2$ by two independent repetitions of themselves we may assume that both have soundness $\tfrac14$ (this affects only constants).
Let $\sigma$ be the encoding of $V_2$ and write $Y$ for the number of symbols of $V_1$'s (extended) proof.

\emph{The verifier.}  $V_{\rm new}$ expects a proof $P_{\rm new}$ of $Y+2^R$ words; the word $P_{\rm new}[j]$ for $j\le Y$ is supposed to be $\sigma(P[j])$, and $P_{\rm new}[Y+w]$ is an auxiliary proof $\pi_w$.
On input $x$ it (1)~chooses $w\in\{0,1\}^R$ and computes $c_w$ and $i_1(w),\dots,i_{Q+1}(w)$ without reading the proof; (2)~chooses a random string $w'$ for $V_2$ and runs $V_2$ on $c_w$ with the
$(Q{+}1)$-part split-encoded assignment $z_j=P_{\rm new}[i_j(w)]$ ($1\le j\le Q+1$) and auxiliary proof $\pi_w=P_{\rm new}[Y+w]$.  It reads only those symbols of the $z_j$ and of $\pi_w$ that $V_2$ queries under $w'$, and accepts iff $V_2$ does.

\emph{Resources.}  The random bits are those of $w$ and $w'$.  Everything else is that of $V_2$ run with $p=Q+1$ parts on a formula of size~$\tau$: alphabet $O(s(\tau))$, queries $O(p+q(\tau))$, decision time $O(p\,t(\tau))$.
The queries depend only on $(w,w')$, so $V_{\rm new}$ is non-adaptive; this is why non-adaptivity was built into every definition.

\emph{Completeness.}  If $x\in L$ take a proof $P$ with $V_1$ accepting on every $w$, and put $P_{\rm new}[j]=\sigma(P[j])$.  By~\eqref{eq:decide} the concatenation of the $P[i_j(w)]$ satisfies $c_w$, so by completeness of $V_2$ there is $\pi_w$ accepted with
probability~$1$; place it at $Y+w$.

\emph{Soundness.}  Let $x\notin L$ and suppose $P_{\rm new}$ is accepted with probability $\ge\tfrac12$.  Define $P[j]=\sigma^{-1}(N(P_{\rm new}[j]))$, where $N$ is the nearest codeword (an arbitrary string if it is undefined).
Let $p_1=\Pr_w[V_1\text{ accepts }P\text{ on }w]<\tfrac14$.  If $V_1$ rejects $P$ on $w$ then by~\eqref{eq:decide} the decoded concatenation $P[i_1(w)]\circ\dots\circ P[i_{Q+1}(w)]$ does not satisfy $c_w$, so
either some $z_j$ is not close to a codeword or the decoded assignment is unsatisfying.  By the soundness of $V_2$, $V_{\rm new}$ accepts on such $w$ with probability $<\tfrac14$, \emph{whatever $\pi_w$ is}.  Hence
\[
  \Pr[V_{\rm new}\text{ accepts}]\ <\ p_1+(1-p_1)\cdot\tfrac14\ <\ \tfrac14+\tfrac34\cdot\tfrac14\ =\ \tfrac7{16}\ <\ \tfrac12,
\]
a contradiction.
\end{proof}

\begin{pitfall}{what composition does \emph{not} do}
Composition does not make a \emph{hard} instance easier; it makes the \emph{queries} cheaper by moving the cost into the proof length and the verifier's computation.
Also, the lemma needs $V_2$ to \emph{check an encoded input}.  A verifier with the same parameters that must read its input in the clear would give nothing: the whole point is that $V_{\rm new}$ never reads a $z_j$ in full.
This is the reason normal form (Section~\ref{sec:normal}), and not just a verifier with small parameters, is the right interface.
\end{pitfall}


%% ===========================================================================
\section{The parameter chase}
\label{sec:chase}
%% ===========================================================================

We now put the pieces together.  Throughout, $n$ is the formula size, $L=\log n$, and all logarithms are base~$2$.

\begin{enumerate}[label=\textbf{Step \arabic*.},leftmargin=*]
  \item \emph{Outer verifier.}  Apply Theorem~\ref{thm:vpoly} with $h+1=2^{\sqrt L}$ and $m=\sqrt L$ ($h\ge\log n$ holds for large $n$).  Call it $V^{(0)}$:
        \[
          (R,S,Q,T)\ =\ \bigl(O(L),\ 2^{O(\sqrt L)},\ \sqrt L,\ 2^{O(\sqrt L)}\bigr).
        \]
  \item \emph{Inner verifier.}  By Section~\ref{sec:normal} (the concatenation test and the Booleanity check added), the same verifier with the same parameter choice is, at input size $N$, an
        $(r,s,q,t)$-constrained normal form verifier with $(r,s,q,t)(N)=(\log N,\ 2^{O(\sqrt{\log N})},\ \sqrt{\log N},\ 2^{O(\sqrt{\log N})})$.
  \item \emph{Compose with $V^{(0)}$ as inner verifier.}  Apply Lemma~\ref{lem:compose}.  Since $T=2^{O(\sqrt L)}$, $\tau=O(T^2)$ has $\log\tau=O(\sqrt L)$, and
        \[
          \begin{aligned}
          \text{random bits}&\;:\; L+O(\log\tau)=O(L),\\
          \text{alphabet bits}&\;:\; 2^{O(\sqrt{\log\tau})}=2^{O(L^{1/4})},\\
          \text{queries}&\;:\; O\bigl(\sqrt L+\sqrt{\log\tau}\bigr)=O(\sqrt L),\\
          \text{decision time}&\;:\; \sqrt L\cdot2^{O(\sqrt{\log\tau})}=2^{O(L^{1/4})}.
          \end{aligned}
        \]
  \item \emph{Iterate.}  Each further composition with $V^{(0)}$ replaces the exponent $e$ in $2^{O(L^{e})}$ by $e/2$ and leaves queries and randomness unchanged up to constant factors.  After $k$ compositions the verifier has
        alphabet bits and decision time $2^{O(L^{2^{-(k+1)}})}$, $O(\sqrt L)$ queries, $O(L)$ random bits.
  \item \emph{Final composition with $V_{\rm savior}$.}  Apply Theorem~\ref{thm:vpoly} with $h=\Theta(\log N)$ and $m=\log N/\log\log N$ at input size $N$: it is a normal form verifier with
        $(r,s,q,t)(N)=(\log N,\ \log N\log\log N,\ \log N/\log\log N,\ \poly\log N)$.  Composing with it, where $\log\tau=O(L^{2^{-(k+1)}})$, gives
        \begin{align*}
          \text{alphabet bits}&=\log\tau\log\log\tau=L^{2^{-(k+1)}+o(1)},\\
          \text{queries}&=O(\sqrt L),\\
          \text{decision time}&=\sqrt L\cdot\poly(\log\tau)=\poly(L).
        \end{align*}
\end{enumerate}

\begin{table}[t]
\centering\small
\renewcommand{\arraystretch}{1.25}
\begin{tabular}{@{}lllll@{}}
\toprule
verifier & random bits & queries (symbols) & bits per symbol & decision time\\
\midrule
$V^{(0)}$ (Step 1) & $O(L)$ & $\sqrt L$ & $2^{O(L^{1/2})}$ & $2^{O(L^{1/2})}$\\
$V^{(1)}=V^{(0)}\circ V^{(0)}$ & $O(L)$ & $O(\sqrt L)$ & $2^{O(L^{1/4})}$ & $2^{O(L^{1/4})}$\\
$V^{(2)}=V^{(1)}\circ V^{(0)}$ & $O(L)$ & $O(\sqrt L)$ & $2^{O(L^{1/8})}$ & $2^{O(L^{1/8})}$\\
$V^{(k)}$ & $O(L)$ & $O(\sqrt L)$ & $2^{O(L^{2^{-(k+1)}})}$ & $2^{O(L^{2^{-(k+1)}})}$\\
$V^{(k)}\circ V_{\rm savior}$ & $O(L)$ & $O(\sqrt L)$ & $L^{2^{-(k+1)}+o(1)}$ & $\poly(L)$\\
\bottomrule
\end{tabular}
\caption{The parameter chase of Section~\ref{sec:chase} ($L=\log n$).  Each composition square-roots the exponent of the alphabet size; the number of queries is stuck at $\sqrt L$ because it is inherited from
the outer verifier.}
\label{tab:chase}
\end{table}

\begin{proof}[Proof of Theorem~\ref{thm:as}]
Given $\eps>0$ choose $k$ with $2^{-(k+1)}\le\eps/2$, a constant.  The verifier of Table~\ref{tab:chase} after $k$ compositions and the final one with $V_{\rm savior}$ uses $O(L)$ random bits,
makes $O(\sqrt L)$ queries to symbols of $L^{\eps/2+o(1)}$ bits each, so reads $L^{1/2+\eps/2+o(1)}\le L^{1/2+\eps}$ bits of the proof, has perfect completeness and soundness $\tfrac12$.
Hence $\sat{3SAT}\in\PCP[\log n,\log^{1/2+\eps}n]$ and so $\NP\subseteq\PCP[\log n,\log^{1/2+\eps}n]$.  The reverse inclusion was noted at the beginning of the lecture.
\end{proof}

\noindent
Each step of the chase hides a constant in $\tau=O(T^2)$ and in $O(Q+q(\tau))$; the next section is about why that matters.

\begin{bigpicture}{What is proved here and what is quoted}
\textbf{Proved in this lecture, from scratch:} Schwartz--Zippel, polynomial extension, the arithmetization (Lemma~\ref{lem:arith}) and zero-testers (Lemma~\ref{lem:zt}), the sum-check (Lemma~\ref{lem:sumcheck}), the exact low-degree
statement (Lemma~\ref{lem:axis}), the concatenation identity (Lemma~\ref{lem:concat}), the composition lemma (Lemma~\ref{lem:compose}), and the whole parameter chase.
\textbf{Quoted:} the two \emph{robust} algebraic lemmas, Lemma~\ref{lem:robust} (consistency of a decomposition) and the robust concatenation theorem (Theorem~4.4.2 of \citealp{AroraSafra98}), both of which rely on
the bivariate lemma of Arora and Safra and on the cubic field size.  Everything else in the text is bookkeeping around them.
\end{bigpicture}


%% ===========================================================================
\section{\texorpdfstring{Consequence: Clique within $2^{\log^{1/2-\eps}N}$ is $\NP$-hard}{Consequence: hardness of approximating Clique}}
\label{sec:clique}
%% ===========================================================================

\begin{proof}[Proof of Corollary~\ref{cor:as-clique}]
Let $0<\eps_1<\eps$.  By Theorem~\ref{thm:as} there is a verifier $V$ for $\sat{3SAT}$ with $O(\log n)$ random bits and $q=\log^{1/2+\eps_1}n$ queries, perfect completeness and soundness $\tfrac12$.  Amplify it by
the random walk on an explicit constant-degree expander on the set of random strings (Exercise~\ref{ex:expander}; \citealp{AKS87}): $k=\lceil\log n/q\rceil$ steps use $O(\log n)+O(k)=O(\log n)$ random bits, make $kq=O(\log n)$ queries, and
give soundness $2^{-\Omega(k)}=2^{-\Omega(\log^{1/2-\eps_1}n)}$.  The FGLSS graph $H_x$ of the amplified verifier (Definition~\ref{def:fglss}) has at most $2^{O(\log n)}=n^{O(1)}=:N$ vertices, and by Theorem~\ref{thm:fglss}
\[
  x\in\sat{3SAT}\Rightarrow\omega(H_x)=2^{r},\qquad x\notin\sat{3SAT}\Rightarrow\omega(H_x)\le2^{r}\cdot2^{-\Omega(\log^{1/2-\eps_1}n)} .
\]
The gap is $2^{\Omega(\log^{1/2-\eps_1}N)}$, which exceeds $2^{c\log^{1/2-\eps}N}$ for every constant $c$ and large $N$ because $\eps_1<\eps$.  An approximation algorithm within $2^{c\log^{1/2-\eps}N}$ would therefore
decide $\sat{3SAT}$ in polynomial time.
\end{proof}

\noindent
Compare with Corollary~\ref{cor:clique-quasi}, which reached $2^{\Omega((\log N)^{1/(1+\eps)})}$ from a constant-query PCP and quasi-polynomial time.  Here the PCP has $\log^{1/2+\eps_1}n$ queries, but the amplification is
randomness-efficient, so the reduction takes \emph{polynomial} time: \emph{this} is the improvement that earned the paper its title.


%% ===========================================================================
\section{\texorpdfstring{Why it stops at $\log^{1/2+\eps}n$, and how to reach $O(1)$}{Why it stops at sublogarithmic query bits, and how to reach O(1)}}
\label{sec:as-limits}
%% ===========================================================================

The chase of Table~\ref{tab:chase} has two bottlenecks.

\paragraph{Queries are inherited, never reduced.}  By Lemma~\ref{lem:compose} the new number of queries is $O(Q+q(\tau))\ge Q$.  Composition shrinks the \emph{symbols} but never the number of symbols queried,
and $V_{\rm poly}$ has $Q=m=\log n/\log(h+1)$.  To have few queries one needs $h$ large, hence symbols of $\Omega(h)$ bits, hence \emph{many} compositions to make them short.

\paragraph{Constants compound.}  With $h+1=n^{1/m}$ and $m$ constant, the alphabet has $n^{\Omega(1/m)}$ bits and shrinking it to $\polylog n$ takes $\Theta(\log\log n)$ compositions, each
multiplying the hidden constants in $O(Q+q(\tau))$ by a factor $c>1$.  After $\Theta(\log\log n)$ rounds the product is $c^{\Theta(\log\log n)}=\polylog n$ queries.  This is exactly the caution in
the paper: ``We do not wish to invoke the Composition Lemma more than $O(1)$ times because it hid constant factors in its $O(\cdot)$ notation and those constant factors could grow very fast.''  The $\log^{1/2+\eps}n$ of Theorem~\ref{thm:as}
is what survives when one limits oneself to $O(1)$ compositions.

\paragraph{The remedy of ALMSS.}  According to Section~6 of \citet{AroraSafra98}, \citet{ALMSS98} proceed in three steps.
\begin{enumerate}
  \item They construct a new outer verifier that reads only $O(1)$ entries of the polynomial tables (not $O(m)$) and has polylogarithmic decision time.  The two ingredients are a more efficient low-degree test, based on the
        main lemma of Arora--Safra (the bivariate Lemma~\ref{lem:robust}'s base case) and on the work of \citet{RubinfeldSudan92}, and a \emph{parallelization} step in the spirit of \citet{LapidotShamir91} and \citet{FeigeLovasz92}
        (a close relative of the constraint--variable game of Lecture~\ref{lec:games}, in which the answers to all queries of a random string are bundled into one symbol), which collapses several queries into one.
  \item They compose this verifier with itself.  The result reads $O(1)$ symbols of size $O(\log\log n)$.
  \item They compose that with a verifier, built from ideas of \citet{BLR90}, which is \emph{wasteful in randomness} (polynomial in the proof size) but makes only $O(1)$ \emph{bit} queries.  Since by now the input
        to the last verifier has size $\polylog n$, its wastefulness is harmless.  The result is a $(\log n,1)$-restricted verifier.
\end{enumerate}
This is Lecture~4.  The last of the three verifiers is based on the \emph{Hadamard code} and linearity testing \citep{BLR90}; the \emph{long code} of \citet{BGS98}, the workhorse of Lecture~7, is a more redundant relative of the Hadamard code.

\begin{bigpicture}{The shape of the argument, in one paragraph}
A satisfying assignment is encoded by a low-degree polynomial, a code in which a $1\%$-corrupted word still determines the codeword.  Satisfiability becomes one vanishing sum, which the sum-check verifies at a single random
point of a polynomial that is itself computed from three random points of the table.  A low-degree test certifies that the table is close to a polynomial.  Because the verifier only ever asks for \emph{uniformly distributed}
points, corruption of a small fraction of the table is invisible to it.  This yields a verifier with $O(\log n)$ random bits and few, but long, queries.  Composition then re-checks each long symbol \emph{in encoded form} with a second such verifier, trading
symbol length for proof length.  Iterating, one reaches sublogarithmic query bits; getting to $O(1)$ needs a better outer verifier and a Hadamard-code inner verifier (Lecture~4), or Dinur's entirely different amplification (Lecture~6).
\end{bigpicture}


%% ===========================================================================
\section*{Exercises}
\addcontentsline{toc}{section}{Exercises}
%% ===========================================================================

\begin{exercise}[\stars{$\star\star$}]\label{ex:sz-tight}
Show that Lemma~\ref{lem:sz} is nearly tight: for $d<q$ and every $m$ there is a nonzero $p\in\mathcal P^m_d$ vanishing on exactly a $1-(1-d/q)^m$ fraction of $\F_q^m$.  Compare $1-(1-d/q)^m$ with $dm/q$.
\end{exercise}

\begin{exercise}[\stars{$\star$}]\label{ex:maj}
Compute the multilinear extension $\hat s$ of the majority function $s\colon\{0,1\}^3\to\{0,1\}$ and evaluate it at $(2,2,2)$ over $\F_7$.  Verify that $\hat s$ has degree $1$ in each variable but total degree $3$.
\end{exercise}

\begin{exercise}[\stars{$\star$}]\label{ex:arith2}
In Example~\ref{ex:arith} take instead $a=(1,1,1,1)$.  Which clause is violated?  Find all nonzero points of $g_{\hat a}$ on $\mathbb K^{8}$ and their values, and compute $\Pr_t[F(t)=0]$ exactly as in Example~\ref{ex:zt}.
\end{exercise}

\begin{exercise}[\stars{$\star\star$}]\label{ex:sc-tight}
Let $B\equiv0$ be the zero polynomial in $\ell$ variables, $H\subseteq\F_q$, $c\ne0$.  Show that the best table convinces the sum-check with probability \emph{exactly} $1-(1-d/q)^\ell$ provided that there is a polynomial
of degree $\le d$ with $d$ distinct roots in $\F_q$ whose sum over $H$ is $c$.  Exhibit one for $q=7$, $d=2$, $H=\{0,1\}$, $c=1$.  (\emph{Hint: scaling a polynomial by a nonzero constant keeps its roots and scales its sum over $H$.})
\end{exercise}

\begin{exercise}[\stars{$\star$}]\label{ex:onedir}
Show that Lemma~\ref{lem:axis} fails if lines of one direction are omitted: for $m=2$ let $f(x_1,x_2)=\varphi(x_1)$ where $\varphi\colon\F_q\to\F_q$ is a function of degree greater than $h$.  Every row $f(x_1,\cdot)$ is constant,
but $f\notin\mathcal P^2_h$.  Prove that $\Delta(f,\mathcal P^2_h)=\Delta(\varphi,\mathcal P^1_h)$, so $f$ can be as far from $\mathcal P^2_h$ as a univariate function can be from $\mathcal P^1_h$.
\emph{Hint: for a fixed polynomial $p$ average over $x_2$.}
\end{exercise}

\begin{exercise}[\stars{$\star\star$}]\label{ex:chase}
Suppose an inner verifier has $s(N)=t(N)=N^{1/3}$, $q(N)=O(1)$, $r(N)=\log N$ and the outer verifier has $(R,S,Q,T)=(O(\log n),n^{1/3},c_0,n^{1/3})$.  Compute the parameters after $k$ compositions
and find the least $k$ for which the alphabet is $\polylog n$.  If each composition multiplies the number of queries by a constant $c>1$, what is the final number of queries?  Compare with Step~4 of Section~\ref{sec:chase}.
\end{exercise}

\begin{exercise}[\stars{$\star\star$}]\label{ex:bool}
Design the Booleanity check mentioned in the pitfalls of Sections~\ref{sec:arith} and~\ref{sec:normal}: give a polynomial in $\mathcal P^{m}_{2h}$ that vanishes on $\mathbb K^m$ iff $A$ is $\{0,1\}$-valued
there, and show how a second sum-check (with its own zero-tester, $k=m$ variables) detects a violation with probability at least $1-3hm/q$.  How many extra queries does the verifier make?
\end{exercise}

\begin{exercise}[\stars{$\star\star\star$}]\label{ex:reading3}
\emph{Reading.}  Compare the low-degree test of Section~4.6 of \citet{AroraSafra98} with the one analyzed in Section~7 of \citet{ALMSS98}: what is tested (axis-parallel lines or random lines), how many queries
are made, how large the field must be, and which of the two lemmas is the bivariate base case.  Which of these differences is responsible for the drop from $O(m)$ to $O(1)$ queries?
\end{exercise}


%% ===========================================================================
\section*{Notes and further reading}
\addcontentsline{toc}{section}{Notes and further reading}
%% ===========================================================================

\paragraph{Where this lecture comes from.}
Everything is from \citet{AroraSafra98}: Sections~3 and~4 give the verifier and the composition lemma, Section~5 the algebraic lemmas that we quoted, and Section~6 the account of what followed.  The algebraic characterization of satisfying polynomials (Lemma~\ref{lem:arith}) is credited there to \citet{BFLS91}, in a form similar in spirit to \citet{BFL91}, and the zero-testers (Lemma~\ref{lem:zt}) to
\citet{BFLS91} and \citet{FGLSS96}; the sum-check is due to \citet{LFKN92}.  The names ``zero-tester'' and ``table of partial sums'', the form of the composition lemma
and our parameter table follow \citet{AroraSafra98}; the numerical examples (Examples~\ref{ex:arith}, \ref{ex:zt}, \ref{ex:sc}) were computed for these notes.

\paragraph{On attributions.}
We distinguish three statements that are often confused: Arora--Safra proved sub-logarithmic queries, Arora--Lund--Motwani--Sudan--Szegedy constant queries.  Section~6 of \citet{AroraSafra98} is generous in crediting the second
paper and the earlier observation of Szegedy and Sudan that a verifier with small decision time yields a gap for \sat{Max-3SAT}.  The modular formulation of composition (PCPs of proximity, assignment testers) is
due to \citet{BGHSV06} and \citet{DR04}; Arora and Safra's normal form verifier is the special case used here.  The analysis of the low-degree test was subsequently improved by \citet{RubinfeldSudan92},
\citet{PolishchukSpielman94} and \citet{AroraSudan97}; none of these is needed for what we do.

\paragraph{A caveat on the text.}
Two small points in the presentation of \citet{AroraSafra98} that we were careful with: the arithmetization certifies rounded satisfiability, so Booleanity must be checked separately when decoding
(pitfall after Lemma~\ref{lem:arith}); and in the remark after Lemma~4.6.4 the paper writes $\gamma^{2(l-1)}\ge\tfrac15$, which is true but too weak for the conclusion drawn from it for $l>5$; the correct bound is $\ge\tfrac45$ (Bernoulli's inequality), and
it is this that makes the choice of $8m/\delta$ checkpoints in the test sufficient.

\paragraph{What is not covered.}
We have not proved the robust lemmas, nor the analysis of $V_{\rm savior}$ at small input size, and we have not treated the $\PCP[\log n,(\log\log n)^2]$ refinement mentioned in the paper.
The proof of the constant-query theorem is Lecture~4, and the combinatorial alternative is Lecture~6.
